
\section{再论\BookRecordIndex{T0256}粘合定理}\label{sec:10-9}

本节考察如下特殊情形下的\BookRecordIndex{T0257}粘合映射：\BookRecordIndex{T0088}殆复结构
$J_z^0=J_z^\infty=:J$ 与 $z\in S^2$ 无关。在这种情形，本章建立的方法给出一个嵌入，
它把一个小二维圆盘与适当\BookRecordIndex{T0191}节点曲线\BookRecordIndex{T0246}模空间的乘积嵌入
$\overline{\mathcal M}_{0,4}(A;J)$。

固定两个同调类 $A^0,A^\infty\in H_2(M;\mathbb Z)$，并记
$A:=A^0+A^\infty$。令 $(T,E,\Lambda)$ 表示带四个标号的树
\[
  T:=(0,\infty),\qquad \Lambda_0:=(1,2),\qquad \Lambda_\infty:=(3,4),
\]
其中只有一条边连接顶点 $0$ 和 $\infty$。设 $J$ 是 $M$ 上被 $\omega$ \BookRecordIndex{T0372}驯服的\BookRecordIndex{T0088}殆复结构。
回顾
\[
  \mathcal M_{0,T}^*(A^{0,\infty};J)
  \subset\mathcal M_{0,T}(A^{0,\infty};J)
\]
表示如下单纯\BookRecordIndex{T0340}稳定映射的\BookRecordIndex{T0246}模空间：它们带四个\BookRecordIndex{T0066}标记点，以 $T$ 为模型，并表示同调类
$A^0,A^\infty$（见第~\ref{sec:5-1} 节和第~\ref{sec:6-1} 节）。把
$\mathcal M_{0,T}(A^{0,\infty};J)$ 的元素写成元组的等价类
\[
  [\mathbf u,\mathbf z]
  =[u^0,z^0,z_1^0,z_2^0;
    u^\infty,z^\infty,z_3^\infty,z_4^\infty],
\]
其中 $u^0,u^\infty:S^2\to M$ 是分别表示同调类 $A^0,A^\infty$ 的
\BookRecordIndex{T0003}$J$-全纯曲线，三元组 $(z^0,z_1^0,z_2^0)$ 和
$(z^\infty,z_3^\infty,z_4^\infty)$ 都由 $S^2$ 上两两不同的点组成，并且
\[
  u^0(z^0)=u^\infty(z^\infty).
\]
等价关系由 \BookRecordIndex{T0040}Möbius 变换偶对群 $G\times G$ 的作用给出；对
$\phi=(\phi^0,\phi^\infty)$，令
\[
  \phi^*(\mathbf u,\mathbf z)
  =\bigl(u^0\circ\phi^0,(\phi^0)^{-1}(z_i^0);
  u^\infty\circ\phi^\infty,(\phi^\infty)^{-1}(z_i^\infty)\bigr).
\]
还要回顾，$\mathcal M_{0,T}(A^{0,\infty};J)$ 是所有带四个\BookRecordIndex{T0066}标记点、表示类 $A$ 的
亏格零\BookRecordIndex{T0340}稳定映射\BookRecordIndex{T0246}模空间 $\overline{\mathcal M}_{0,4}(A;J)$ 的子集。记
\[
  \iota_T:\mathcal M_{0,T}(A^{0,\infty};J)
  \longrightarrow\overline{\mathcal M}_{0,4}(A;J)
\]
为自然包含映射。

对每个 $[\mathbf u,\mathbf z]\in\mathcal M_{0,T}(A^{0,\infty};J)$，存在唯一的
\BookRecordIndex{T0040}Möbius 变换偶对 $\phi^0,\phi^\infty\in G$，使得
\begin{equation}\label{eq:10-9-1}
\begin{gathered}
  \phi^0(z^0)=0,\qquad \phi^0(z_1^0)=1,\qquad \phi^0(z_2^0)=\infty,\\
  \phi^\infty(z^\infty)=\infty,\qquad
  \phi^\infty(z_3^\infty)=1,\qquad
  \phi^\infty(z_4^\infty)=0.
\end{gathered}
\end{equation}
对 $c>0$，定义
\[
  \mathcal M_{0,T}^*(c):=
  \left\{[\mathbf u,\mathbf z]\in\mathcal M_{0,T}^*(A^{0,\infty};J)
  \ \middle|\
  (u^0\circ\phi^0,u^\infty\circ\phi^\infty)\in\mathcal M^*(c)
  \right\},
\]
其中 $\phi^0,\phi^\infty$ 取为满足式~\eqref{eq:10-9-1} 的变换，而
$\mathcal M^*(c)\subset\mathcal M^*(A^{0,\infty};J)$ 是由条件 (6.2.5) 和
(6.2.6) 定义的紧集穷竭族。对 $\rho>0$，令 $B_\rho\subset\mathbb C$ 为以原点为
中心、半径为 $\rho$ 的闭圆盘。

\begin{theorem}\label{thm:10-9-1}
设 $J$ 是 $M$ 上被 $\omega$ \BookRecordIndex{T0372}驯服的\BookRecordIndex{T0088}殆复结构，并且在定义~\ref{def:6-2-1} 的意义下
对 $(A^0,A^\infty)$ 正则。对每个常数 $c>0$，存在常数 $\rho>0$ 和连续嵌入
\[
  \iota_{T,c}:\mathcal M_{0,T}^*(c)\times B_\rho
  \longrightarrow\overline{\mathcal M}_{0,4}(A;J),
\]
具有下列性质。
\begin{enumerate}[label=(\roman*)]
\item $\iota_{T,c}$ 在 $\mathcal M_{0,T}^*(c)\times\{0\}$ 上的限制是自然包含
$\iota_T$。
\item $\iota_{T,c}$ 的像是
$\overline{\mathcal M}_{0,4}(A;J)$ 中
$\iota_T(\mathcal M_{0,T}^*(c-1))$ 的一个邻域。
\item $\iota_{T,c}$ 在流形
$\mathcal M_{0,T}^*(c)\times(B_\rho\setminus\{0\})$ 上的限制，是到
$\mathcal M_{0,4}(A;J)$ 中的光滑保定向嵌入。
\item 对 $i=1,2,3,4$，令
\[
  \operatorname{ev}_i:\overline{\mathcal M}_{0,4}(A;J)\longrightarrow M
\]
表示在第 $i$ 个\BookRecordIndex{T0066}标记点处的\BookRecordIndex{T0283}求值映射。对 $\lambda\in B_\rho$，定义
$\operatorname{ev}_{i,\lambda}:\mathcal M_{0,T}^*(c)\to M$ 为复合
\[
  \operatorname{ev}_{i,\lambda}([\mathbf u,\mathbf z])
  :=\operatorname{ev}_i\circ
  \iota_{T,c}([\mathbf u,\mathbf z],\lambda).
\]
则当 $\lambda$ 趋于零时，
$\operatorname{ev}_{i,\lambda}:\mathcal M_{0,T}^*(c)\to M$ 在 $C^1$ 拓扑下收敛到
$\operatorname{ev}_i\circ\iota_T:\mathcal M_{0,T}^*(c)\to M$。
\end{enumerate}
\end{theorem}

\begin{remark}\label{rem:10-9-2}
由于 $J$ 不依赖于 $z$，一般而言，只能通过一般地选取 $J$ 来使单纯\BookRecordIndex{T0340}稳定映射达到正则。
这就是定理~\ref{thm:10-9-1} 援引定义~\ref{def:6-2-1} 而非定义~\ref{def:10-1-1}
的原因，也是定理~\ref{thm:10-1-2} 的\BookRecordIndex{T0257}粘合映射在这种情形只定义于子集
$\mathcal M^*(c)\subset\mathcal M(c)$、而不是整个 $\mathcal M(c)$ 上的原因。
\end{remark}

\begin{proof}[定理~\ref{thm:10-9-1} 的证明]
令 $0<\delta_0<1$ 如定理~\ref{thm:10-1-2} 中所取，并固定
$0<\delta<\delta_0$。对 $R>1/(\delta\delta_0)$，用
\[
  \iota_c^{\delta,R}:\mathcal M^*(c)\longrightarrow\mathcal M(A;J)
\]
表示定理~\ref{thm:10-1-2} 的\BookRecordIndex{T0257}粘合映射。任取满足下式的正数 $\rho$
\sourcefootnote{原书取 $\rho=(\delta\delta_0)^2$，同时把 $B_\rho$ 定义为闭圆盘，
却在随后对其所有非零元素使用严格不等式 $|\lambda|<\rho$，从而不能覆盖边界。
定理只要求某个正半径，故改为任取 $0<\rho<(\delta\delta_0)^2$；于是闭圆盘上的
$R=|\lambda|^{-1/2}$ 也始终满足所需的严格下界。}：
\[
  0<\rho<(\delta\delta_0)^2
\]
并如下定义
\[
  \iota_{T,c}:\mathcal M_{0,T}^*(c)\times B_\rho
  \longrightarrow\overline{\mathcal M}_{0,4}(A;J).
\]
当 $\lambda=0$ 时，定义
$\iota_{T,c}([\mathbf u,\mathbf z],0):=[\mathbf u,\mathbf z]$。现在固定
$\lambda\in B_\rho\setminus\{0\}$，并写成
\[
  \lambda=\frac{e^{i\theta}}{R^2}.
\]
由 $|\lambda|\leq\rho<(\delta\delta_0)^2$，有
$R>1/(\delta\delta_0)$。$\mathcal M_{0,T}^*(c)$ 中的每个等价类都有唯一代表元
$(\mathbf u,\mathbf z)$，满足
\begin{equation}\label{eq:10-9-2}
  z^0=0,\quad z_1^0=1,\quad z_2^0=\infty,\qquad
  z^\infty=\infty,\quad z_3^\infty=e^{i\theta},\quad z_4^\infty=0.
\end{equation}
该代表元满足 $u^0(0)=u^\infty(\infty)$，并且
$(u^0,u^\infty)\in\mathcal M^*(c)$。最后一个条件成立，是因为空间
$\mathcal M^*(c)$ 在 $S^1\times S^1$ 的重参数化作用下不变。定义如下元素
\sourcefootnote{原书在这里定义元素时漏写输入 $\lambda$，印作
$\iota_{T,c}([\mathbf u,\mathbf z])$；本映射的定义域是
$\mathcal M_{0,T}^*(c)\times B_\rho$，而右端也显式依赖 $\lambda$，故补全。}：
\[
  \iota_{T,c}([\mathbf u,\mathbf z],\lambda)
  =[u_\lambda,\mathbf z_\lambda]
  \in\overline{\mathcal M}_{0,4}(A;J)
\]
其中
\[
  u_\lambda:=\iota_c^{\delta,R}(u^0,u^\infty),
  \qquad \mathbf z_\lambda:=(1,\infty,\lambda,0).
\]
映射
\[
  \mathcal M_{0,T}^*(c)\times(B_\rho\setminus\{0\})
  \longrightarrow\mathcal M_{0,4}(A;J),
  \qquad ([\mathbf u,\mathbf z],\lambda)\longmapsto
  [u_\lambda,\mathbf z_\lambda]
\]
良定义，因为每个等价类只有一个满足式~\eqref{eq:10-9-2} 的代表元。此外，它依定义
满足条件（i），并且在
$\mathcal M_{0,T}^*(c)\times(B_\rho\setminus\{0\})$ 上光滑，因为由
定理~\ref{thm:10-1-2} 第（i）部分，映射
\[
  (1/(\delta\delta_0),\infty)\times\mathcal M^*(c)
  \longrightarrow\mathcal M(A;J),
  \qquad(R,u^0,u^\infty)\longmapsto
  \iota_c^{\delta,R}(u^0,u^\infty)
\]
是光滑的。\sourcefootnote{原书在这个光滑映射的定义域中把固定 $J$ 情形所需的
$\mathcal M^*(c)$ 漏掉星号，写成 $\mathcal M(c)$；同页单射性论证又把
$\mathcal M_{0,T}^*(c)$ 写成未定义的 $\mathcal M_{0,T}(c)$。由
注~\ref{rem:10-9-2} 及本证明开头\BookRecordIndex{T0257}粘合映射的定义域，两处均补上星号。}

下面证明
\[
  \iota_{T,c}:\mathcal M_{0,T}^*(c)\times B_\rho
  \longrightarrow\overline{\mathcal M}_{0,4}(A;J)
\]
是连续嵌入。当 $\lambda\ne0$ 时，点 $([\mathbf u,\mathbf z],\lambda)$ 处的连续性
来自它在 $\mathcal M_{0,T}^*(c)\times(B_\rho\setminus\{0\})$ 上的光滑性；点
$([\mathbf u,\mathbf z],0)$ 处的连续性来自定理~\ref{thm:10-1-2} 第（ii）部分的收敛
结论。为说明该映射单射，设
$([\mathbf u,\mathbf z],\lambda)$ 和
$([\mathbf u',\mathbf z'],\lambda')$ 是
$\mathcal M_{0,T}^*(c)\times B_\rho$ 的两个元素，且
\[
  \iota_{T,c}([\mathbf u,\mathbf z],\lambda)
  =\iota_{T,c}([\mathbf u',\mathbf z'],\lambda').
\]
于是 $\mathbf z_\lambda=\mathbf z_{\lambda'}$，从而 $\lambda=\lambda'$；再由
定理~\ref{thm:10-1-2} 中映射
$\iota_c^{\delta,R}:\mathcal M^*(c)\to\mathcal M(A;J)$ 的单射性，有
$[\mathbf u,\mathbf z]=[\mathbf u',\mathbf z']$。所以 $\iota_{T,c}$ 是连续单射。
它是恰当的，因为其定义域紧。因此 $\iota_{T,c}$ 是连续嵌入。

接着证明，当 $\rho$ 充分小时，$\iota_{T,c}$ 的像是
$\iota_T(\mathcal M_{0,T}^*(c-1))$ 的一个邻域。否则，存在不在其像中的序列
\[
  [u^\nu,\mathbf z^\nu]\in\overline{\mathcal M}_{0,4}(A;J),
\]
\BookRecordIndex{T0141}格罗莫夫收敛到某个
$[\mathbf u,\mathbf z]\in\mathcal M_{0,T}^*(c-1)$。于是
$[u^\nu,\mathbf z^\nu]$ 不可能属于\BookRecordIndex{T0072}层
$\mathcal M_{0,T}^*(A^{0,\infty};J)$，因而必属于
$\mathcal M_{0,4}(A;J)$。不妨把 $[u^\nu,\mathbf z^\nu]$ 的参数化选成满足
\[
  z_1^\nu=1,\qquad z_2^\nu=\infty,\qquad
  z_3^\nu=: \lambda^\nu,\qquad z_4^\nu=0.
\]
由于 $[u^\nu,\mathbf z^\nu]$ \BookRecordIndex{T0141}格罗莫夫收敛到
$\mathcal M_{0,T}^*(c-1)$ 中的元素，序列 $\lambda^\nu$ 必收敛到零。因此
$R^\nu:=|\lambda^\nu|^{-1/2}\to\infty$；必要时取子列，可以假设
$\arg\lambda^\nu=: \theta^\nu$ 收敛到 $\theta$。还可以不妨假设
$[\mathbf u,\mathbf z]$ 对这个 $\theta$ 满足归一化条件~\eqref{eq:10-9-2}。

序列 $[u^\nu,\mathbf z^\nu]$ \BookRecordIndex{T0141}格罗莫夫收敛到 $[\mathbf u,\mathbf z]$，意味着
$u^\nu$ 在 $S^2\setminus\{0\}$ 上\BookRecordIndex{T0416}逐紧一致收敛到 $u^0$，并且存在 \BookRecordIndex{T0040}Möbius 变换序列
$\phi^\nu$，使 $u^\nu\circ\phi^\nu$ 在 $S^2\setminus\{\infty\}$ 上\BookRecordIndex{T0416}逐紧一致
收敛到 $u^\infty$，且
\[
  \infty=\lim_{\nu\to\infty}(\phi^\nu)^{-1}(\infty),
  \qquad e^{i\theta}=\lim_{\nu\to\infty}(\phi^\nu)^{-1}(\lambda^\nu),
  \qquad 0=\lim_{\nu\to\infty}(\phi^\nu)^{-1}(0).
\]
把 $\phi^\nu$ 与另一个收敛到恒等映射的 \BookRecordIndex{T0040}Möbius 变换序列复合后，可以假设
\[
  \phi^\nu(z)=z/(R^\nu)^2,
  \qquad z\in S^2=\mathbb C\cup\{\infty\}.
\]
所以当 $\nu$ 很大时，映射 $u^\nu$ 对 $R=R^\nu$ 和偶对
$(u^0,u^\infty)\in\mathcal M^*(c-1)$ 满足定理~\ref{thm:10-1-2} 第（v）部分的
条件。这蕴含存在
$(\widetilde u^0,\widetilde u^\infty)\in\mathcal M^*(c)$，使得
\[
  u^\nu=\iota_c^{\delta,R^\nu}(\widetilde u^0,\widetilde u^\infty).
\]
因此依构造，
\[
  [u^\nu,\mathbf z^\nu]
  =\iota_{T,c}
  ([\widetilde u^0,0,1,\infty;
     \widetilde u^\infty,\infty,e^{i\theta^\nu},0],\lambda^\nu).
\]
这与假设矛盾，从而证明定理~\ref{thm:10-9-1} 第（ii）部分。第（iii）、（iv）部分
直接来自定理~\ref{thm:10-1-2} 的\BookRecordIndex{T0257}粘合映射 $\iota_c^{\delta,R}$ 的相应性质。
定理~\ref{thm:10-9-1} 得证。
\end{proof}

定理~\ref{thm:10-9-1} 的证明中构造的\BookRecordIndex{T0257}粘合映射还具有下述性质。\BookRecordIndex{T0255}粘合参数
\begin{equation}\label{eq:10-9-3}
  \lambda=w(z_4,z_1,z_2,z_3)
  =\frac1{w(z_1,z_2,z_3,z_4)}
\end{equation}
恰是\BookRecordIndex{T0257}粘合映射像中元素
$\iota_{T,c}([\mathbf u,\mathbf z],\lambda)$ 的\BookRecordIndex{T0066}标记点\BookRecordIndex{T0185}交比。在附加这个条件后，
殆复情形下不能指望\BookRecordIndex{T0283}求值映射
\[
  \operatorname{ev}_i\circ\iota_{T,c}:
  \mathcal M_{0,T}^*(c)\times B_\rho\longrightarrow M
\]
连续可微。不过，通过额外的重标度可以做到连续可微。

\begin{corollary}\label{cor:10-9-3}
定理~\ref{thm:10-9-1} 中的\BookRecordIndex{T0257}粘合映射
\[
  \iota_{T,c}:\mathcal M_{0,T}^*(c)\times B_\rho
  \longrightarrow\overline{\mathcal M}_{0,4}(A;J)
\]
可选成使\BookRecordIndex{T0283}求值映射
\[
  \operatorname{ev}_i\circ\iota_{T,c}:
  \mathcal M_{0,T}^*(c)\times B_\rho\longrightarrow M
\]
对 $i=1,2,3,4$ 都连续可微。
\end{corollary}

\begin{proof}
选取光滑函数 $\beta:[0,\infty)\to[0,\infty)$，使
$\dot\beta(0)=0$，且当 $s>0$ 时 $\dot\beta(s)>0$。把
定理~\ref{thm:10-9-1} 证明中的映射 $\iota_{T,c}$ 换成下列映射
\sourcefootnote{原书在这个定义的右端再次写成 $\iota_{T,c,\beta}$，使公式成为自指；
这里是在原\BookRecordIndex{T0257}粘合映射 $\iota_{T,c}$ 的参数中作额外重标度，故补正右端下标。}：
\[
  \iota_{T,c,\beta}([\mathbf u,\mathbf z],\lambda)
  :=\iota_{T,c}([\mathbf u,\mathbf z],\beta(|\lambda|)\lambda).
\]
由定理~\ref{thm:10-9-1} 第（iv）部分，
$\operatorname{ev}_i\circ\iota_{T,c,\beta}$ 关于
$\mathcal M_{0,T}^*(c)$ 中变量的导数，在
$\mathcal M_{0,T}^*(c)\times B_\rho$ 上连续。关于变量
$\lambda\in B_\rho$ 的偏导数连续延拓到
$\mathcal M_{0,T}^*(c)\times B_\rho$，并且在
$\mathcal M_{0,T}^*(c)\times\{0\}$ 上消失。最后一步来自
定理~\ref{thm:10-1-2} 证明中\BookRecordIndex{T0257}粘合映射的构造；细节留作读者的习题。
这证明了推论~\ref{cor:10-9-3}。
\end{proof}

本书中，\BookRecordIndex{T0254}粘合的主要应用是证明\BookRecordIndex{T0143}格罗莫夫--威滕不变量的\BookRecordIndex{T0125}分裂公理（见
定理~\ref{thm:7-5-10}），这已在第~\ref{sec:10-8} 节完成。那里讨论的一般情形相当
复杂。不过，我们对\BookRecordIndex{T0125}分裂公理的主要应用是证明量子乘法的结合律，而为此只需要
$\#I=4$ 时的\BookRecordIndex{T0125}分裂公理；比较命题 11.1.11。在这种情形，只须知道：在\BookRecordIndex{T0185}交比
$\lambda$ 固定的一组点处满足四个约束的曲线计数，等于曲线偶对计数的加权和；偶对中
每条曲线都通过两个约束以及一个适当的基类 $e_\nu$。上面定理~\ref{thm:10-9-1} 的表述
使这种对应变得清楚。它蕴含：当 $|\lambda|>0$ 充分小时，\BookRecordIndex{T0257}粘合映射
$\iota_{T,c}$ 的逆在前一种计数与穿过给定约束的相交曲线偶对计数之间建立双射。

为完成证明，还须使用“分裂对角线”的标准技巧。换言之，要用到 $M\times M$ 中的对角线
庞加莱对偶于
\[
  \sum_{\nu,\mu}g^{\nu\mu}\pi_1^*(e_\nu)\smile\pi_2^*(e_\mu)
\]
这一事实；比较式~\eqref{eq:10-8-2}。

\begin{exercise}\label{ex:10-9-4}
用定理~\ref{thm:10-9-1} 代替定理~\ref{thm:10-1-2}，并按需使用
定理~\ref{thm:10-8-1} 的有关细节，完成这一论证。
\end{exercise}
